\documentclass[../../../include/open-logic-section]{subfiles}

\begin{document}

\olfileid{sth}{ord-arithmetic}{mult}

\olsection{الضرب الترتيبي}

ننتقل الآن إلى الضرب الترتيبي، ونتناوله على نحو قريب جدًا من تناولنا
الجمع الترتيبي. افترض إذن أننا نريد ضرب~$\alpha$ في~$\beta$. ويمكنك
لفعل ذلك أن تتخيل شبكة مستطيلة عرضها~$\alpha$ وارتفاعها~$\beta$؛
فيكون حاصل ضرب~$\alpha$ و$\beta$ هو نتيجة التحرك عبر كل صف، ثم
الانتقال عبر الصف التالي\ldots{} إلى أن تجتاز الشبكة كلها. وبعبارة
أخرى، ينتج حاصل ضرب~$\alpha$ و$\beta$ من إحلال نسخة من~$\alpha$ محل
\emph{كل} عنصر في~$\beta$.

ولصياغة ذلك صوريًا، نستخدم ببساطة الترتيب المعجمي المعكوس على حاصل
الضرب الديكارتي لـ$\alpha$ و$\beta$:

\begin{defn}
لأي عددين ترتيبيين~$\alpha, \beta$، يكون حاصل ضربهما
$\alpha \ordtimes \beta = \ordtype{\alpha \times \beta, \rlexless}$.
\end{defn}
\noindent
ينبغي لنا مرة أخرى أن نتحقق من أن هذا تعريف حسن الصياغة:

\begin{lem}\ollabel{ordtimeslessiswo}
$\tuple{\alpha \times \beta, \rlexless}$ ترتيب جيد، لأي عددين
ترتيبيين~$\alpha$ و$\beta$.
\end{lem}

\begin{proof}
تمامًا كما في~\olref[add]{ordsumlessiswo}.
\end{proof}
\noindent
وليس من العسير إثبات أن الضرب يسلك هكذا:

\begin{lem}\ollabel{ordtimesrecursion}
لأي عددين ترتيبيين~$\alpha, \beta$:
\begin{align*}
	\alpha \ordtimes 0 &= 0\\
	\alpha \ordtimes (\beta \ordplus 1) &=
		(\alpha \ordtimes \beta) \ordplus \alpha\\
	\alpha  \ordtimes \beta &=
		\supstrict_{\delta < \beta}(\alpha \ordtimes \delta) &&
		\text{عندما يكون $\beta$ عددًا ترتيبيًا حديًا}.
\end{align*}
\end{lem}

\begin{proof}
متروك تمرينًا.
\end{proof}

والواقع أنه، مثلما في حالة الجمع، كان بوسعنا تعريف الضرب الترتيبي
بهذه المعادلات العودية بدلًا من تقديم تعريف مباشر. وبالمثل أيضًا،
كما في الجمع، ثمة سلوك مألوف:

\begin{lem}\ollabel{ordinalmultiplicationisnice}
إذا كانت~$\alpha, \beta, \gamma$ أعدادًا ترتيبية، فإن:
\begin{enumerate}
	\item\ollabel{ordtimes1} إذا كان~$\alpha \neq 0$ و$\beta < \gamma$،
	فإن~$\alpha \ordtimes \beta < \alpha \ordtimes \gamma$؛
	\item\ollabel{ordtimes2} إذا كان~$\alpha \neq 0$ و
	$\alpha \ordtimes \beta = \alpha\ordtimes\gamma$، فإن
	$\beta = \gamma$؛
	\item\ollabel{ordtimes3} $\alpha \ordtimes (\beta \ordtimes
	\gamma) = (\alpha \ordtimes \beta) \ordtimes \gamma$؛
	\item\ollabel{ordtimes4} إذا كان~$\alpha \leq \beta$، فإن
	$\alpha \ordtimes \gamma \leq \beta \ordtimes \gamma$؛
	\item\ollabel{ordtimes5} $\alpha \ordtimes (\beta \ordplus
	\gamma) = (\alpha \ordtimes \beta )\ordplus (\alpha\ordtimes
	\gamma)$.
\end{enumerate}
\end{lem}

\begin{proof}
متروك تمرينًا.
\end{proof}

يمكنك إثبات نتائج أخرى (أو البحث عنها) بقدر ما تشاء. لكن، في ضوء
\olref[ord-arithmetic][add]{ordsumnotcommute}، ينبغي ألا تكون النتيجة
الآتية مفاجئة:

\begin{prop}
الضرب الترتيبي ليس تبادليًا:
$2 \ordtimes \omega = \omega < \omega \ordtimes 2$
\end{prop}

\begin{proof}
$2 \ordtimes \omega = \supstrict_{n < \omega}(2\ordtimes n) = \omega
\in \supstrict_{n < \omega}(\omega \ordplus n) = \omega \ordplus
\omega = \omega \ordtimes 2$.
\end{proof}
\noindent
ومرة أخرى، المسوغ الحدسي مباشر جدًا. فلكي تحسب
$2 \ordtimes \omega$، تستبدل بكل عدد طبيعي كيانين. وستحصل على نمط
الترتيب نفسه إذا أدخلت ببساطة جميع الأعداد «النصفية» بين الأعداد
الطبيعية، أي إذا نظرت في الترتيب الطبيعي على
$\Setabs{\nicefrac{n}{2}}{n \in \omega}$. وعند عرض الأمر بهذه الطريقة،
يكون نمط الترتيب هو بوضوح نمط~$\omega$ نفسه. أما لكي تحسب
$\omega \ordtimes 2$، فإنك تضع نسختين من~$\omega$، إحداهما بعد
الأخرى.

\begin{prob}
أثبت
\olref[sth][ord-arithmetic][mult]{ordtimeslessiswo}،
و\olref[sth][ord-arithmetic][mult]{ordtimesrecursion}،
و
\olref[sth][ord-arithmetic][mult]{ordinalmultiplicationisnice}
\end{prob}

\end{document}
